\documentclass[10pt,a4paper]{amsart}
\usepackage[margin=19mm]{geometry}
\usepackage{amsmath,amssymb,mathtools}
\usepackage[scr=rsfs,cal=euler]{mathalfa}
\usepackage{xfrac}
\usepackage[unicode,hidelinks,bookmarks=false]{hyperref}
\newcommand{\ie}{i.e.\ }
\newcommand{\cf}{cf.\ }
\newcommand{\resp}{respectively\ }
\newcommand{\Subsection}{\subsection}
\providecommand{\nobreakdash}{\nobreak\hbox{-}}
\setlength{\emergencystretch}{2em}
\multlinegap0pt
\usepackage{bm}
\usepackage{bbm}
\usepackage{dsfont}
\usepackage{xspace}
\usepackage{cancel}
\usepackage{mathtools}
\usepackage{amsthm}
\makeatletter
\@ifundefined{thm}{\newtheorem{thm}{Thm}}{}
\@ifundefined{lem}{\newtheorem{lem}[thm]{Lemma}}{}
\makeatother
\newcommand{\bbT}{\mathbb{T}}
\newcommand{\bx}{\mathbf{x}}
\newcommand{\dv}{\operatorname{dv}}
\newcommand{\id}{\operatorname{id}}
\newcommand{\sbse}{\subseteq}
\title[Minimizing point configurations for tensor product energies ]{Minimizing point configurations for tensor product energies on the torus}
\author{Dmitriy Bilyk, Nicolas Nagel, Ian Ruohoniemi}
\date{Source: arXiv:2510.25442v1 (CC BY 4.0)}
\begin{document}
\maketitle

\noindent\emph{Source and licence.} Minimizing point configurations for tensor product energies on the torus. Version arXiv:2510.25442v1 (CC BY 4.0), distributed under the Creative Commons Attribution 4.0 International licence, as recorded in the arXiv licence metadata for this exact version and shown on the abstract page https://arxiv.org/abs/2510.25442v1. Full rights notice: licenses/arxiv-2510.25442-CC-BY-4.0.txt.\par\smallskip

\noindent\textit{Editorial scope.} Counted unit: the Lem whose source label is \texttt{lem:double\_counting} in the article (source0.tex lines 543--551, source label \texttt{lem:double\_counting}), delivered together with the article's own complete original proof of it (source0.tex lines 553--594, one \texttt{proof} environment of 42 lines). No mathematics is written, repaired or re-derived: statement and proof are the article's own text line for line. Required context retained uncounted: none. Every definition, display and label of those lines is retained unchanged. Omitted from this package and not redistributed: 1--542, 552--552, 595--2594. Nothing inside a retained excerpt is omitted or altered: every retained body is delivered exactly as the article's lines are written, so no omission receipt of any kind accompanies this package. Citations: the retained excerpts carry no citation command at all, so no bibliography entry of the article is transcribed and no reference is redistributed. Reference adaptation: none was needed and none was made; every reference inside the delivered unit resolves inside it, so no unresolved reference or citation is produced. Printed slips kept literal: no printed word, symbol, hypothesis or number of the counted statement or of its proof was altered. Layout and class: the article's own class is \texttt{article} with packages including hyperref, fontenc, inputenc, babel, mathtools, amsfonts, color, amsmath, amsthm, amssymb, mathrsfs, amstext, float, todonotes; the wrapper is the lane's pinned \texttt{amsart} editorial wrapper at 10pt on A4, which restates the theorem-like environments the retained text needs and adds the article's own macro definitions that the retained text needs (\texttt{\textbackslash bbT}, \texttt{\textbackslash bx}, \texttt{\textbackslash dv}, \texttt{\textbackslash id}, \texttt{\textbackslash sbse}), with extra wrapper packages (bm, bbm, dsfont, xspace, cancel, mathtools). The delivered numbering is the unit's own. Assets: no retained span contains a figure environment, a graphics-inclusion command, a PDF-page inclusion, a table or a TikZ picture; no image file is copied into this package, the package declares no asset dependency and no image file is redistributed with it. Licence: version arXiv:2510.25442v1 of this article is distributed under the Creative Commons Attribution 4.0 International licence, as recorded in the arXiv licence metadata for this exact version and shown on the abstract page https://arxiv.org/abs/2510.25442v1. A full rights notice is delivered at licenses/arxiv-2510.25442-CC-BY-4.0.txt. Characters: every retained excerpt is pure ASCII, so the delivered engine, XeLaTeX with its default Latin Modern text font, reports no missing character.
\begin{lem} \label{lem:double_counting}
    Let $X \sbse \bbT^d, \#X=N$ be a Latin hypercube set with a  single distance vector, whose common distance vector is given by $\{(\rho_1, \dots, \rho_d)\} = \dv(X)$. Then for every $r = 1, \dots, \lfloor N/2 \rfloor$ it holds
    $$
    \#\left\{i=1, \dots, d: \rho_i = \frac rN\right\} =  \begin{cases}
        \frac{2d}{N-1}, & r < \frac N2, \\
        \frac d{N-1}, &  r = \frac N2.
    \end{cases}
    $$
\end{lem}

\begin{proof}
    We will show, via double counting, that
    \begin{align*}
        & (N^2-N) \cdot \#\left\{i=1, \dots, d: \rho_i = \frac rN\right\} \\
        = & d \cdot \#\left\{(k, \ell) \in \{0, 1, \dots, N-1\}^2: \delta\left(\frac kN, \frac \ell N\right) = \frac rN\right\},
    \end{align*}
    from which the claim follows since
    $$
    \#\left\{(k, \ell) \in \{0, 1, \dots, N-1\}^2: \delta\left(\frac kN, \frac \ell N\right) = \frac rN\right\} = \begin{cases}
        2N, &  r < \frac N2, \\
        N ,&  r = \frac N2.
    \end{cases}
    $$
    %By Definition \ref{def:lat_hypcube} 
    Since $X$ is a Latin hypercube, there are permutations $\sigma_1, \dots, \sigma_d$ %of $\{0, 1, \dots, N-1\}$ 
    such that % with
    $$
    X = \{\bx^k\}_{k=0}^{N-1} = \left\{\left(\frac{\sigma_1(k)}N, \dots, \frac{\sigma_d(k)}N\right): k=0, 1, \dots, N-1\right\}
    $$
    (with $\sigma_1 = \id$ the identity on $\{0, 1, \dots, N-1\}$). Consider the array $A = [a_{(k, \ell), i}]$ with rows indexed by $(k, \ell) \in \{0, 1, \dots, N-1\}^2, k \neq \ell$ and columns indexed by $i=1, \dots, d$, with entries
    $$
    a_{(k, l), i} = \delta\left(\frac{\sigma_i(k)}N, \frac{\sigma_i(\ell)}N\right).
    $$
    We will count for how many entries it holds $a_{(k, \ell), i} = \frac rN$. The row corresponding to $(k, \ell)$ consists, by definition, of the entries of $\dv(\bx^k, \bx^\ell)$ in some order. By the single distance property, this row contains exactly
    $$
    \#\left\{i=1, \dots, d: \rho_i = \frac rN\right\}
    $$
    such entries, so that there are
    $$
    (N^2-N) \cdot \#\left\{i=1, \dots, d: \rho_i = \frac rN\right\}
    $$
    in total of them in $A$. On the other hand, since $\sigma_i$ is a permutation, in the $i$-th column we have
    \begin{align*}
        & \#\left\{(k, \ell) \in \{0, 1, \dots, N-1\}^2: \delta\left(\frac {\sigma_i(k)}N, \frac {\sigma_i(\ell)}N\right) = \frac rN\right\} \\
        = & \#\left\{(k, \ell) \in \{0, 1, \dots, N-1\}^2: \delta\left(\frac {k}N, \frac {\ell}N\right) = \frac rN\right\},
    \end{align*}
    so in all of $A$ there are
    $$
    d \cdot \#\left\{(k, \ell) \in \{0, 1, \dots, N-1\}^2: \delta\left(\frac {k}N, \frac {\ell}N\right) = \frac rN\right\}
    $$
    such entries. This finishes the proof.
\end{proof}

\end{document}
